\chapter{Complexity — Path Walk}
%%% generated by the blueprint pipeline from TCSlib.Complexity.SpaceComplexity.Logspace.PathWalk
\section{Overview}

Declarations of the Lean module \texttt{TCSlib.Complexity.SpaceComplexity.Logspace.PathWalk}.

\section{Declarations}

\begin{definition}[Walk relation of bounded length]
\lean{Complexity.WalkLe}
\statementsource{Vad12}{pdf-bf2ce5422088-p036-b001}
\label{Complexity.WalkLe}
\leanok
Let $E$ be a relation on natural numbers and let $n$ be a natural number. Define
$\operatorname{WalkLe}(j,c,t)$ recursively: for $j=0$, it means $c=t$; for $j+1$, it
means either $c=t$, or there is a natural number $v<n$ such that $E(c,v)$ and
$\operatorname{WalkLe}(j,v,t)$.
\end{definition}

\begin{lemma}[Reflexive bounded walks]
\lean{Complexity.walkLe_refl}
\label{Complexity.walkLe_refl}
\leanok
\uses{Complexity.WalkLe}
\difficulty{3}
For every relation $E$ on $\mathbb{N}$ and all $n,j,c\in\mathbb{N}$, there is a walk
from $c$ to itself of length at most $j$, where a walk of length at most zero means
equality and a walk at a positive bound means either equality or an initial step to some
$v<n$ following $E$, followed by a walk with bound one less.
\end{lemma}

\begin{lemma}[Monotonicity of bounded walks]
\lean{Complexity.walkLe_succ}
\label{Complexity.walkLe_succ}
\leanok
\uses{Complexity.WalkLe}
\difficulty{4}
Let $E$ be a relation on the natural numbers and let $n$ be a natural number. If $c$
reaches $t$ by a walk of length at most $j$ that uses only vertices less than $n$ (with
a walk of length zero meaning $c=t$), then $c$ reaches $t$ by a walk of length at most
$j+1$ under the same vertex restriction.
\end{lemma}

\begin{lemma}[Monotonicity of bounded walks]
\lean{Complexity.walkLe_mono}
\label{Complexity.walkLe_mono}
\leanok
\uses{Complexity.WalkLe, Complexity.walkLe_succ}
\difficulty{4}
Let $E$ be a binary relation on the natural numbers and let $n,j,k,c,t\in\mathbb{N}$. If
$j\le k$ and there is a walk from $c$ to $t$ of length at most $j$, with every
intermediate vertex less than $n$ and each step related by $E$, then there is such a
walk of length at most $k$. A walk of length at most $j$ here means equality when $j=0$,
and for a positive bound either equality or an initial step to some $v<n$ with $E(c,v)$
followed by a walk of length at most $j-1$ from $v$ to $t$.
\end{lemma}

\begin{lemma}[Bounded walks give restricted reachability]
\lean{Complexity.walkLe_reach}
\label{Complexity.walkLe_reach}
\leanok
\uses{Complexity.WalkLe}
\difficulty{4}
Let $E$ be a relation on the natural numbers, and let $n,j,c,t\in\mathbb{N}$ with $c<n$.
If there is a walk from $c$ to $t$ of length at most $j$, in which each initial vertex
of a step is less than $n$ and each step follows $E$, then $t$ is reachable from $c$ by
a finite sequence of steps whose endpoints are both less than $n$ and whose steps follow
$E$. A walk of length at most $j$ is defined recursively: at bound zero it means $c=t$,
and at a positive bound it means either $c=t$ or there is $v<n$ with $E(c,v)$ and a walk
of length at most $j-1$ from $v$ to $t$.
\end{lemma}

\begin{lemma}[Finite bound for a reflexive transitive walk]
\lean{Complexity.exists_walkLe}
\label{Complexity.exists_walkLe}
\leanok
\uses{Complexity.WalkLe}
\difficulty{4}
Let $E$ be a binary relation on the natural numbers and let $n\in\mathbb{N}$. If $t$ is
reachable from $c$ by a finite sequence of steps, each from some $u<n$ to some $v<n$
with $E(u,v)$, then there is a natural number $j$ such that $t$ can be reached from $c$
using at most $j$ steps of this form.
\end{lemma}

\begin{lemma}[Shortening the walk bound]
\lean{Complexity.walkLe_short}
\label{Complexity.walkLe_short}
\leanok
\uses{Complexity.WalkLe, Complexity.walkLe_mono}
\difficulty{6}
Let $E$ be a relation on the natural numbers and let $n,j,c,t\in\mathbb{N}$, with $c<n$
and $t<n$. If there is a walk from $c$ to $t$ of length at most $j$, then there is one
of length at most $n$. A walk of length at most $j$ is defined recursively: at bound
zero it means equality, and at a positive bound it means either equality or an initial
step to some $v<n$ following $E$, followed by a walk with bound one less.
\end{lemma}

\begin{lemma}[Reachability by a bounded walk]
\lean{Complexity.reach_iff_walkLe}
\label{Complexity.reach_iff_walkLe}
\leanok
\uses{Complexity.WalkLe, Complexity.exists_walkLe, Complexity.walkLe_reach, Complexity.walkLe_short}
\difficulty{4}
Let $E$ be a relation on the natural numbers, let $n\in\mathbb{N}$, and let $c,t<n$.
There is a finite path from $c$ to $t$ using steps of $E$ whose endpoints are both less
than $n$ if and only if there is a walk from $c$ to $t$ of length at most $n$, with
every intermediate vertex less than $n$. A walk of length zero is allowed.
\end{lemma}

\begin{lemma}[List Index Bound]
\lean{Complexity.lt_of_getD_true}
\label{Complexity.lt_of_getD_true}
\leanok
\difficulty{3}
For a Boolean list $l$ and a natural number $i$, if the value of $l$ at index $i$, using
false as the default when the index is out of range, is true, then $i<|l|$.
\end{lemma}

\begin{lemma}[Extensionality for equal-length Boolean lists]
\lean{Complexity.ext_getD'}
\label{Complexity.ext_getD'}
\leanok
\difficulty{3}
If two Boolean lists have the same length and their entries agree at every index smaller
than that length (using false as the default value for out-of-range indices), then the
lists are equal.
\end{lemma}

\begin{definition}[Doubled word at a position]
\lean{Complexity.DblAt}
\label{Complexity.DblAt}
\leanok
For Boolean words $x,u$ and a natural number $q$, say that $u$ is doubled at position
$q$ in $x$ if, for every $i<|u|$, the entries of $x$ at positions $q+2i$ and $q+2i+1$
both equal the entry of $u$ at position $i$. The entries of $x$ at positions $q+2|u|$
and $q+2|u|+1$ must respectively be false and true; any out-of-range list lookup is
interpreted as false.
\end{definition}

\begin{definition}[Word at a position]
\lean{Complexity.PlainAt}
\label{Complexity.PlainAt}
\leanok
For Boolean words $x,u$ and a natural number $r$, say that $u$ occurs at position $r$ in
$x$ if, for every $i<|u|$, the entry of $x$ at position $r+i$ equals the entry of $u$ at
position $i$, and $|x|=r+|u|$. Any out-of-range list lookup is interpreted as false.
\end{definition}

\begin{lemma}[Suffix after pair encoding]
\lean{Complexity.drop_pairEncode}
\label{Complexity.drop_pairEncode}
\leanok
\uses{Turing.pairEncode, Turing.pairEncode_eq_dbl}
\difficulty{3}
For any two finite bit strings $u$ and $w$, deleting the first $2|u|+2$ bits from their
pair encoding leaves exactly $w$.
\end{lemma}

\begin{lemma}[Characterization by the Doubled Pair Encoding]
\lean{Complexity.dblAt_iff}
\label{Complexity.dblAt_iff}
\leanok
\uses{Complexity.DblAt, Complexity.ext_getD', Complexity.getD_drop, Complexity.lt_of_getD_true, Complexity.pairEncode_getD, Turing.length_pairEncode, Turing.pairEncode}
\difficulty{5}
For every pair of binary words $x,u$ and every natural number $q$, the condition that
$x$ has the doubled encoding of $u$ starting at position $q$—each bit of $u$ repeated
twice, followed by $01$—holds if and only if the suffix of $x$ beginning at $q$ is the
concatenation of that doubled encoding and the suffix of $x$ beginning at $q+2|u|+2$.
\end{lemma}

\begin{lemma}[Characterization of a plain segment]
\lean{Complexity.plainAt_iff}
\label{Complexity.plainAt_iff}
\leanok
\uses{Complexity.PlainAt, Complexity.ext_getD', Complexity.getD_drop}
\difficulty{4}
For Boolean words $x$ and $u$ and a natural number $r$, $x$ agrees with $u$ starting at
position $r$ and has length $r+|u|$ if and only if $r\le |x|$ and deleting the first $r$
entries of $x$ gives $u$.
\end{lemma}

\begin{definition}[Raw path encoding]
\lean{Complexity.pathRaw}
\statementsource{Sav70}{
  pdf-c8033d2fd705-p012-b002,
  pdf-c8033d2fd705-p012-b004
}
\label{Complexity.pathRaw}
\leanok
\uses{Turing.pairEncode}
For a natural number $n$, a Boolean word $M$, and natural numbers $s,t$, define the raw
path encoding by pairing the length-$n$ all-true word with the result of pairing $M$
with the encoding obtained by pairing the binary digit lists of $s$ and $t$.
\end{definition}

\begin{definition}[Raw adjacency relation]
\lean{Complexity.rawEdge}
\statementsource{Sav70}{
  pdf-c8033d2fd705-p012-b002,
  pdf-c8033d2fd705-p012-b003
}
\label{Complexity.rawEdge}
\leanok
For a natural number $n$, a Boolean list $M$, and natural numbers $u,v$, declare $u$
adjacent to $v$ when $u<n$, $v<n$, and the entry of $M$ at index $un+v$ is true (an
out-of-range entry is treated as false).
\end{definition}

\begin{definition}[Raw path language]
\lean{Complexity.pathRawLang}
\statementsource{Vad12}{pdf-bf2ce5422088-p035-b005}
\statementsource{Sav70}{
  pdf-c8033d2fd705-p012-b003,
  pdf-c8033d2fd705-p012-b004,
  pdf-c8033d2fd705-p012-b005
}
\label{Complexity.pathRawLang}
\leanok
\uses{Complexity.pathRaw, Complexity.rawEdge}
The language consists of Boolean lists $x$ for which there are natural numbers $n,s,t$
and a Boolean list $M$ of length $n^2$ such that $s,t<n$, $x$ is the path encoding
determined by $n,M,s,t$, and there is a finite sequence of raw-adjacency steps from $s$
to $t$.
\end{definition}

\begin{lemma}[Characterization of the raw path encoding]
\lean{Complexity.eq_pathRaw_iff}
\label{Complexity.eq_pathRaw_iff}
\leanok
\uses{Complexity.DblAt, Complexity.PlainAt, Complexity.dblAt_iff, Complexity.drop_pairEncode, Complexity.pathRaw, Complexity.plainAt_iff, Turing.length_pairEncode, Turing.pairEncode}
\difficulty{5}
Let $x$ and $M$ be finite Boolean words and let $n,s,t$ be natural numbers. Then $x$ is
the encoding of the pair consisting of $n$ true bits and the nested pair of $M$ with the
pair of the binary expansions of $s$ and $t$ (using the encoding that doubles each bit
of the first component, appends $01$, and then appends the second component) if and only
if all of the following hold: starting at position $0$, the first $n$ bits are each
doubled and followed by the separator $01$; starting at position $2n+2$, the bits of $M$
are each doubled and followed by $01$; starting at position $2n+2+2|M|+2$, the bits of
the binary expansion of $s$ are each doubled and followed by $01$; and starting at
position $2n+2+2|M|+2+2|\text{binary expansion of }s|+2$, the binary expansion of $t$
appears verbatim and ends exactly at the end of $x$. In each doubled segment, both
copies of every bit agree with the corresponding bit of that segment's word; positions
beyond the end of $x$ are interpreted as false when checking these conditions.
\end{lemma}

\begin{lemma}[Row-major index bound]
\lean{Complexity.rowMajor_lt}
\label{Complexity.rowMajor_lt}
\leanok
\difficulty{2}
Let $n,u,v\in\mathbb{N}$ with $u<n$ and $v<n$. Then $un+v<n^2$.
\end{lemma}

\begin{definition}[Encoded path edge]
\lean{Complexity.pathEdge}
\label{Complexity.pathEdge}
\leanok
For a Boolean list $x$ and natural numbers $n,u,v$, the encoded path relation holds when
the entry of $x$ at index $2n+2+2(un+v)$ is true, with out-of-range entries treated as
false.
\end{definition}

\begin{definition}[Unary block terminator]
\lean{Complexity.UnaryRest}
\label{Complexity.UnaryRest}
\leanok
For a Boolean list $x$ and natural numbers $k,n$, every $i$ with $k\le i<n$ has true
entries at indices $2i$ and $2i+1$; the entries at indices $2n$ and $2n+1$ are
respectively false and true. Out-of-range entries are treated as false.
\end{definition}

\begin{definition}[Matrix block terminator]
\lean{Complexity.MatRest}
\label{Complexity.MatRest}
\leanok
For a Boolean list $x$ and natural numbers $n,i$, each $j$ with $i\le j<n^2$ has equal
entries at indices $2n+2+2j$ and $2n+2+2j+1$. The entries at indices $2n+2+2n^2$ and
$2n+2+2n^2+1$ are respectively false and true; out-of-range entries are treated as
false.
\end{definition}

\begin{definition}[Path acceptance condition]
\lean{Complexity.PathAcc}
\statementsource{Sav70}{
  pdf-c8033d2fd705-p012-b002,
  pdf-c8033d2fd705-p012-b003,
  pdf-c8033d2fd705-p012-b004,
  pdf-c8033d2fd705-p012-b005
}
\label{Complexity.PathAcc}
\leanok
\uses{Complexity.DblAt, Complexity.MatRest, Complexity.PlainAt, Complexity.UnaryRest, Complexity.WalkLe, Complexity.pathEdge}
A Boolean list $x$ satisfies the acceptance condition if there exist natural numbers
$n,s,t$ such that the unary block condition holds with parameters $0,n$, the matrix
block condition holds with parameters $n,0$, $s<n$, and the doubled encoding of the
binary digits of $s$ begins at index $2n+2+2n^2+2$. Also, $t<n$, the plain encoding of
the binary digits of $t$ begins at index $2n+2+2n^2+2+2\ell+2$, where $\ell$ is the
number of binary digits of $s$, and there is a walk from $s$ to $t$ of length at most
$n$ using the encoded path relation on an $n$-by-$n$ vertex set.
\end{definition}

\begin{lemma}[Unary doubled prefix characterization]
\lean{Complexity.unaryRest_iff}
\label{Complexity.unaryRest_iff}
\leanok
\uses{Complexity.DblAt, Complexity.UnaryRest}
\difficulty{3}
For every binary word $x$ and $n\in\mathbb N$, the following are equivalent: at each
pair of positions $2i,2i+1$ for $0\le i<n$, both bits of $x$ are $1$, and the next pair
is $01$; or the first $n$ bits of the word encoded by those doubled positions are all
$1$, with the separator $01$ immediately afterward. In all indexed-bit statements, a
position beyond the end of $x$ is read as $0$.
\end{lemma}

\begin{lemma}[Characterization of the matrix remainder condition]
\lean{Complexity.matRest_iff}
\label{Complexity.matRest_iff}
\leanok
\uses{Complexity.DblAt, Complexity.MatRest}
\difficulty{5}
For every Boolean word $x$ and natural number $n$, the following are equivalent. First,
for every $j<n^2$, the two bits of $x$ at positions $2n+2+2j$ and $2n+3+2j$ agree, and
the bits at positions $2n+2+2n^2$ and $2n+3+2n^2$ are respectively $0$ and $1$. Second,
there is a Boolean word $M$ of length $n^2$ such that, for each $i<n^2$, the bits of $x$
at positions $2n+2+2i$ and $2n+3+2i$ both equal the $i$th bit of $M$, and the same final
two positions contain $0$ and $1$, respectively. Out-of-range word positions are
interpreted as $0$.
\end{lemma}

\begin{theorem}[Characterization of the path encoding]
\lean{Complexity.pathAcc_iff}
\label{Complexity.pathAcc_iff}
\leanok
\uses{Complexity.DblAt, Complexity.PathAcc, Complexity.eq_pathRaw_iff, Complexity.matRest_iff, Complexity.pathEdge, Complexity.pathRawLang, Complexity.reach_iff_walkLe, Complexity.rowMajor_lt, Complexity.unaryRest_iff}
\difficulty{6}
For every bit string $x$, the following are equivalent. First, there are $n\in\mathbb
N$, an $n\times n$ Boolean adjacency matrix $M$, and vertices $s,t<n$ such that $x$
begins with $n$ pairs $11$, followed by $01$; then each matrix bit is represented by an
equal pair of bits, followed by $01$; then the binary expansion of $s$ is represented by
equal pairs followed by $01$, and the binary expansion of $t$ follows verbatim and ends
the string. In addition, there is a walk from $s$ to $t$ using at most $n$ edges of the
graph encoded by $M$. Second, $x$ is of the form consisting of the doubled unary
encoding of $n$, followed by the similarly encoded matrix $M$, then the similarly
encoded binary expansion of $s$, then that of $t$, and $t$ is reachable from $s$ in the
directed graph whose vertices are $0,\ldots,n-1$ and whose edges are the true entries of
$M$.
\end{theorem}

\begin{lemma}[Least significant bit decomposition]
\lean{Complexity.bits_eq_cons}
\label{Complexity.bits_eq_cons}
\leanok
\uses{Complexity.LogProg.bits_half, Complexity.LogProg.bits_head_odd}
\difficulty{4}
For every nonzero natural number $w$, its binary expansion written from least
significant bit to most significant bit begins with the bit indicating whether $w$ is
odd, followed by the binary expansion of $\lfloor w/2\rfloor$.
\end{lemma}

\begin{lemma}[Initial Bit Decomposition of a Doubled Encoding]
\lean{Complexity.dblAt_cons}
\label{Complexity.dblAt_cons}
\leanok
\uses{Complexity.DblAt}
\difficulty{5}
For every pair of binary words $x,u$, every natural number $q$, and every bit $b$, the
condition that $x$ has the doubled encoding of $b::u$ starting at position $q$—meaning
each bit of $b::u$ occurs twice at its prescribed positions and is followed by the
delimiter $01$—holds if and only if $x$ has $b$ at positions $q$ and $q+1$ and has the
doubled encoding of $u$ starting at position $q+2$.
\end{lemma}

\begin{lemma}[Plain prefix after its first bit]
\lean{Complexity.plainAt_cons}
\label{Complexity.plainAt_cons}
\leanok
\uses{Complexity.PlainAt}
\difficulty{4}
For Boolean words $x$ and $u$, a natural number $r$, and a bit $b$, $x$ agrees with the
word $b::u$ starting at position $r$ and has length $r+|b::u|$ if and only if $x_r=b$
and $x$ agrees with $u$ starting at position $r+1$ and has length $(r+1)+|u|$. Here
out-of-range entries are read as false.
\end{lemma}

\begin{lemma}[Quotient and remainder of a successor]
\lean{Complexity.divmod_succ}
\label{Complexity.divmod_succ}
\leanok
\difficulty{3}
Let $n>0$ and $i$ be natural numbers. If the remainder of $i$ upon division by $n$,
increased by one, equals $n$, then the quotient of $i+1$ upon division by $n$ is the
quotient of $i$ plus one, and the new remainder is zero. If the remainder of $i$
increased by one is not $n$, then the quotient is unchanged and the new remainder is the
previous remainder plus one.
\end{lemma}

\begin{lemma}[Characterization of a quotient equal to the divisor]
\lean{Complexity.div_eq_self_iff}
\label{Complexity.div_eq_self_iff}
\leanok
\difficulty{4}
For natural numbers $n$ and $i$ with $i\le n^2$, one has $i/n=n$ if and only if $i=n^2$.
\end{lemma}
